\scanpage{34}
Il faut de plus des conditions qui assurent,
à partir des données (18), que $i\mapsto X_i$
est croissant, et \deleted{que} la condition (14)
$X_i\cap X_j=\bigcup_{k\le i,j}X_k$. Pour \deleted{les} exprimer ces conditions,
on introduit pour tt $\beta\in I_0$ l'ens
\[
(21)\qquad I_\beta^*=\{i\in I^*\mid i>\beta\}
\]
donc la donnée de $(I\supset\unsure{I'|I_0})$
\inserted{($I$ ens. ordonné)} équivaut à celle
de $I^*,I_0$ \inserted{(ens. ordonné \quad ens.)} et de la famille des
$I_\beta^*\subset I^*$ ($\beta\in I_0$).

Pour \deleted{vérifier}\inserted{exprimer} que \deleted{si}
$i\le j$ \inserted{(donc $j\in I^*$)} implique $X_i\subset X_j$,
on distingue les cas où $i\in I_0$\inserted{$=\beta$}
\deleted{\unsure{ou}}, où c'est
clair sur (19), et celui où $i\in I^*$ où \deleted{et} c'est
clair aussi sur (\unsure{19|17}), compte tenu de la condition
néc. évidente
\[
\text{e)}\qquad \deleted{\unsure{\forall}} i,j\in I^*,\ i\le j
\quad\Longrightarrow\quad X_i^*\subset X_j^*
\]
\deleted{Pour que}\inserted{On doit} avoir aussi, manifestement
\[
\text{f)}\qquad \unsure{j|i},j\in I^*\Longrightarrow
X_i^*\cap X_j^*=\bigcup_{k\le i,j}X_k^*
\]

Pour la relation
\[
X_i\cap X_j=\bigcup_{k\le i,j}X_k\qquad(i,j\in I)
\]
il \deleted{de} suffit de prendre le cas où $i,j$ ne
sont pas comparables. Quand un des deux
est $\in I_0$, alors cette relation \deleted{est} signifie que
$\beta\in I_0$, $i\in I^*\setminus I_\beta$, on a
$X_\beta\cap X_j=\varnothing$, ce qui résulte de (19).
\margin{33}
